\chapter{2021 Nobel Prize in Economics}
\textbf{Laureates: }David Card (USA/Canada), Joshua Angrist (USA), Guido Imbens (USA/Netherlands)

\section*{Reason for Award}
For their methodological contributions to the analysis of causal relationships

\section{What They Did}
David Card, Joshua Angrist, and Guido Imbens advanced the empirical analysis
of causal relationships, providing tools for drawing credible causal inferences
from observational data.

Card studied the minimum wage, showing that moderate increases did not reduce
employment as standard theory predicted. His research on the New Jersey-Pennsylvania
minimum wage increase challenged conventional wisdom and demonstrated that
labor markets do not always behave as simple supply-demand models predict.
Card's natural experiment approach became a model for applied microeconomics.

Angrist and Imbens developed the local average treatment effect (LATE) framework,
showing how instrumental variables identify causal effects for subpopulations.
They demonstrated that instrumental variable estimates identify the effect for
compliers---those whose treatment status is affected by the instrument---rather
than the average treatment effect for the entire population.

Their work clarified when and how natural experiments, regression discontinuity
designs, and difference-in-differences methods can identify causal effects.
These methods allow researchers to estimate causal impacts without randomized
experiments, which are often impossible or unethical.

Card's research on immigration showed that increased immigration had minimal
effects on native workers' employment and wages. His work on education examined
the returns to schooling and class size effects. Angrist's research on the
Vietnam draft lottery estimated the causal effect of military service on
earnings.

\section{Impact on Economic Thinking}
Card, Angrist, and Imbens transformed empirical economics by establishing
credibility revolutions in causal inference.

Their work showed that observational data could yield credible causal estimates
when combined with creative research designs. This influenced policy evaluation
across education, labor, health, and development economics.

The local average treatment effect framework clarified interpretation of
instrumental variable estimates, preventing overgeneralization of results.
Researchers learned to be precise about which population their estimates apply to.

Card's minimum wage research changed how economists think about labor markets,
showing that empirical evidence can contradict theoretical predictions. The
empirical tradition of testing theory against data became more prominent.

Card's research on immigration showed that increased immigration had minimal
effects on native workers' employment and wages. This finding challenged
political rhetoric and provided evidence-based guidance for immigration policy.

Angrist's work on the Vietnam draft lottery used a natural experiment to
estimate the causal effect of military service on earnings. The random
assignment of draft numbers provided exogenous variation in military service.

Their methodological contributions have been widely adopted, making empirical
economics more rigorous and policy-relevant. The emphasis on identification
strategies has raised standards for causal inference across social sciences.
Fields such as political science, sociology, and public health have adopted
these methods.

\section{Modern Perspectives Today}
The methodological advances of Card, Angrist, and Imbens remain central to
empirical research. Natural experiments, regression discontinuity designs,
and difference-in-differences are standard tools in economics and related
fields.

Research on education policy, labor markets, health economics, and development
evaluates interventions using these methods. The credibility revolution has
become the dominant paradigm in applied microeconomics.

The COVID-19 pandemic generated numerous causal studies using these approaches
to assess policy impacts. Randomized evaluations of lockdown effects, stimulus
payments, and vaccination programs built on the methodological foundations
established by the laureates.

Recent work extends their methods to heterogeneous treatment effects, dynamic
settings, and machine learning applications. Research on causal inference in
high-dimensional settings and with big data continues to build on their
foundational contributions.

The credibility revolution continues to influence how economists design studies,
evaluate evidence, and make policy recommendations. Their work demonstrates the
power of creative research designs to identify causal relationships in complex
real-world settings.

\section*{Important Bibliography}
\begin{itemize}
\item Card, D., \& Krueger, A. (1994). ``Minimum wages and employment: A case study of the fast-food industry in New Jersey and Pennsylvania.'' \textit{American Economic Review}, 84(4), 772--793.
\item Angrist, J. (1990). ``Lifetime earnings and the Vietnam era draft lottery: Evidence from social security administrative records.'' \textit{American Economic Review}, 80(3), 313--336.
\item Angrist, J., \& Imbens, G. (1994). ``Identification of causal effects using instrumental variables.'' \textit{Journal of the American Statistical Association}, 91(434), 444--455.
\item Imbens, G., \& Rubin, D. (2015). \textit{Causal Inference for Statistics, Social, and Biomedical Sciences}. Cambridge University Press.
\item Angrist, J., \& Pischke, J. (2009). \textit{Mostly Harmless Econometrics: An Empiricist's Companion}. Princeton University Press.
\end{itemize}
